%%%%%%%%%%%%%%%%%%%%%%%%%%%%%%%%%%%%%%%%%
% Academic Title Page
% LaTeX Template
% Version 2.0 (17/7/17)
%
% This template was downloaded from:
% http://www.LaTeXTemplates.com
%
% Original author:
% WikiBooks (LaTeX - Title Creation) with modifications by:
% Vel (vel@latextemplates.com)
%
% License:
% CC BY-NC-SA 3.0 (http://creativecommons.org/licenses/by-nc-sa/3.0/)
% 
% Instructions for using this template:
% This title page is capable of being compiled as is. This is not useful for 
% including it in another document. To do this, you have two options: 
%
% 1) Copy/paste everything between \begin{document} and \end{document} 
% starting at \begin{titlepage} and paste this into another LaTeX file where you 
% want your title page.
% OR
% 2) Remove everything outside the \begin{titlepage} and \end{titlepage}, rename
% this file and move it to the same directory as the LaTeX file you wish to add it to. 
% Then add \input{./<new filename>.tex} to your LaTeX file where you want your
% title page.
%
%%%%%%%%%%%%%%%%%%%%%%%%%%%%%%%%%%%%%%%%%

%----------------------------------------------------------------------------------------
%	PACKAGES AND OTHER DOCUMENT CONFIGURATIONS
%----------------------------------------------------------------------------------------

\documentclass[11pt]{article}

\usepackage[utf8]{inputenc} % Required for inputting international characters
\usepackage[T1]{fontenc} % Output font encoding for international characters
\usepackage{graphicx}
\usepackage{mathpazo} % Palatino font
\usepackage[legalpaper,margin=0.9in] {geometry}
\setcounter{tocdepth}{3}
\setcounter{secnumdepth}{3}

\begin{document}

%----------------------------------------------------------------------------------------
%	TITLE PAGE
%----------------------------------------------------------------------------------------

\begin{titlepage} % Suppresses displaying the page number on the title page and the subsequent page counts as page 1
	\newcommand{\HRule}{\rule{\linewidth}{0.5mm}} % Defines a new command for horizontal lines, change thickness here
	
	\center % Centre everything on the page
	
	%------------------------------------------------
	%	Headings
	%------------------------------------------------
	
	\textsc{\LARGE University of Regina}\\[1.5cm] % Main heading such as the name of your university/college
	
	\textsc{\Large Software Systems Engineering}\\[0.5cm] % Major heading such as course name
	
	
	
	%------------------------------------------------
	%	Title
	%------------------------------------------------
	
	\HRule\\[0.4cm]
	
	{\Huge\bfseries Workshop Enterprise Resource Planning Suite Resources and Specifications Document}\\[0.4cm] % Title of your document
	
	\HRule\\[1.5cm]
	
	%------------------------------------------------
	%	Author(s)
	%------------------------------------------------
	
	\begin{minipage}{0.4\textwidth}
		\begin{flushleft}
			\large
			Jonathan Wells\\
			\textsc{200328640} % Your name
		\end{flushleft}
	\end{minipage}
	~
	\begin{minipage}{0.4\textwidth}
		\begin{flushright}
			\large
			Konstantin Kharitonov\\
			\textsc{200354502} % Supervisor's name
		\end{flushright}
	\end{minipage}
	
	% If you don't want a supervisor, uncomment the two lines below and comment the code above
	%{\large\textit{Author}}\\
	%John \textsc{Smith} % Your name
	%------------------------------------------------
	%	Logo
	%------------------------------------------------
	
	\vfill\vfill
	\includegraphics[width=0.7\textwidth]{UR.png}\\[2cm] % Include a department/university logo - this will require the graphicx package
	 

	%------------------------------------------------
	%	Date
	%------------------------------------------------
	
	\vfill\vfill\vfill % Position the date 3/4 down the remaining page
	
	{\large\today} % Date, change the \today to a set date if you want to be precise
	
	%----------------------------------------------------------------------------------------
	
	\vfill % Push the date up 1/4 of the remaining page
	
\end{titlepage}

%----------------------------------------------------------------------------------------

%----------------------------------------------------------------------------------------
%Table of Contents %

\newpage 
\tableofcontents
%-------------------------------------------------------------------------------------

\newpage

%-------------------------------------------------------------------------------------
%Introduction%
\section{Introduction}
The Workshop Enterprise Resource Planning Suite, or ERP for short, is an administrative task management web application primarily designed for the Engineering Workshop at University of Regina main campus. It is to be the main application to be used for managing incoming and outgoing workorders, which are student and faculty submitted forms requesting the service of the shop. The service provides workorder capacity planning, allowing for the user to actively manage the status of each project, time tracking features for large scale and small scale work, as well as the ability to track the inventory of the workshop, including but not limited to, materials, tools and equipment. 
\subsection{Purpose}
This system was designed to replace the previous methods of workorder, time, and inventory tracking, centralizing all aspects into one powerful application that can be accessed online. Workorders currently must be submitted via paper form directly to the workshop during its operating hours. The form must then be reviewed by the workshop manager and if accepted, future meetings are scheduled. All workorders submitted are then stored physically in binders, which date back to the opening of the workshop. All materials and inventory are also stored physically. This project intends to automate all workorders and have then be submitted and archived electronically. As well, the system is intended to track all scopes of projects, ranging from small miscellaneous tasks to larger scale projects in such a fashion that the workshop manager can schedule them effectively in advance. 
\subsection{Scope}
ERP is designed as a Web API, such that it is run in browser and is able to be accessed from any computer with a sufficient internet connection. It will be a local application that will be primarily accessed by the workshop manager, who is this project's main client. Secondary clients include faculty and staff that wish to submit workorders over the ERP suite. The primary client is the only one intended to have full control of all features of the ERP suite. 
\newline
{\setlength{\parindent}{0cm}

The ERP Suite currently is planned to be exclusive to the engineering workshop based on its design as of the completion of this capstone project, as future work on this project will require a redesign to be re-purposed for future clients. The ideal future client for this program is for machine and workshop owners with a staff less than 50.  

\section{Overall Description}
\section{System Features}
\section{External Interface Requirements}
\section{Non-Functional Requirements}
\section{References}


\end{document}
